% =============================================================================
%  Section 5 — PCA + Logistic Regression
% =============================================================================
\section{Principal Component Analysis}
\label{sec:acp}

PCA is applied separately to the three splits, on the CT~$-$~T
\emph{deltas} of the economic and HLTV features, after standard
scaling (\texttt{StandardScaler}). We pick \texttt{StandardScaler}
over a robust variant because the upstream cleaning steps (corrupted
demos, phantom rounds, 4v5 pistols) have already removed the
structural outliers. The aim is to test whether an effective
dimensionality reduction is possible before the regression.

\begin{table}[H]
\centering
\small
\begin{tabular}{lc}
\toprule
\textbf{Regime} & \textbf{Axes for 80\,\% inertia / initial features} \\
\midrule
\texttt{pistol}      & 12 / 70  \\
\texttt{post-pistol} & 12 / 105 \\
\texttt{normal}      & 10 / 106 \\
\bottomrule
\end{tabular}
\caption{PCA summary by regime.}
\label{tab:pca_summary}
\end{table}

None of the three regimes lends itself to a useful reduction: 10 to
12 axes are needed to capture 80\,\% of cumulative inertia, i.e. about
15\,\% of the initial features. We therefore keep the full feature set
for the logistic regression.

% =============================================================================
\section{Logistic regression}
\label{sec:logreg}

For each regime, we model the probability that CT win the round
through a logistic regression:
\[
P(Y=1 \mid X = x) \;=\; \sigma\!\left(\alpha + \beta^{\!\top} x\right),
\qquad \sigma(t) = \frac{1}{1+e^{-t}},
\]
where $Y = 1$ indicates a CT win. The train/test split is
\textbf{temporal and grouped by match}: matches are sorted by
\texttt{match\_date} and the most recent 20\,\% form the test set
(80/20). This scheme preserves match integrity (no round from the
same match is split between train and test) and reproduces the
conditions of a real prediction setup (past $\rightarrow$ future).
Features are standardised on the training set only. We compare a
baseline (majority class) to four nested feature configurations:
\textbf{A} (round-state only), \textbf{B} = A + HLTV,
\textbf{C} = B + map (one-hot), and a case \textbf{D} that adds
quasi-target variables (\texttt{rw\_pct}, \texttt{pistol\_win\_pct},
round-2 indicators) in order to measure the effect of including
variables that already partly encode the target. We report
\textbf{accuracy} and \textbf{AUC} on the test set.

\begin{table}[H]
\centering
\small
\renewcommand{\arraystretch}{1.05}
\begin{tabular}{lccc}
\toprule
\textbf{AUC} \scriptsize{(accuracy in parentheses)} & \textbf{Pistol} & \textbf{Normal} & \textbf{Post-pistol} \\
\midrule
Baseline (majority class)     & 0.50~\scriptsize{(.50)} & 0.50~\scriptsize{(.51)} & 0.50~\scriptsize{(.53)} \\
A — round-state               & \textbf{0.54}~\scriptsize{(.54)} & \textbf{0.70}~\scriptsize{(.63)} & \textbf{0.86}~\scriptsize{(.78)} \\
B — A + HLTV                  & \textbf{0.55}~\scriptsize{(.54)} & \textbf{0.71}~\scriptsize{(.64)} & \textbf{0.85}~\scriptsize{(.76)} \\
C — A + HLTV + map            & \textbf{0.56}~\scriptsize{(.51)} & \textbf{0.71}~\scriptsize{(.64)} & \textbf{0.85}~\scriptsize{(.76)} \\
D — + quasi-target features   & \textit{0.56}~\scriptsize{(.54)} & \textit{0.71}~\scriptsize{(.64)} & \textit{0.85}~\scriptsize{(.76)} \\
\bottomrule
\end{tabular}
\caption{AUC and accuracy on the test set, by regime and feature
configuration.}
\label{tab:logreg}
\end{table}

Three levels of predictability emerge:
\[
\text{pistol}\;\text{AUC}\!\approx\!0.56 \;<\;
\text{normal}\;\text{AUC}\!\approx\!0.71 \;<\;
\text{post-pistol}\;\text{AUC}\!\approx\!0.86.
\]
The pistol regime stays barely above chance, in line with the absence
of a strong individual signal observed in
Section~\ref{sec:split_pearson}. On the other two regimes, the
\emph{round state} alone (case A) already captures most of the
predictive power. The HLTV, map and even quasi-target additions
\textbf{add little} (at most $+0.01$ AUC on the normal and post-pistol
regimes): the carrying signal is concentrated in the immediate
economic and material state, and the global team indicators only
complement it marginally.

% =============================================================================
